%------------------------------------------------------------------------------%

\usepackage{integracao_e_series}

\title{Teorema Fundamental do Cálculo}

%------------------------------------------------------------------------------%
\begin{document}

\addtitle

%------------------------------------------------------------------------------%
\section{Teorema do Valor Médio}

%------------------------------------------------------------------------------%
\begin{frame}{Derivadas e Integrais}
\emph{Derivada} de $f$: \quad taxa de variação de $f$
\[
  f'(x) = \frac{df}{dx}
\]

\pause
\emph{Integral indefinida} de $f$: \quad antiderivada de $f$
\[
  F(x) = \int f(x)\,dx \qquad\Rightarrow\qquad \frac{dF}{dx} = f(x)
\]

\pause
\emph{Integral definida} de $f$: \quad diferença de áreas sob o gráfico de $f$
\[
  \text{``Área''} = \int_a^b f(x)\,dx
\]
\end{frame}

%------------------------------------------------------------------------------%
\begin{frame}{Teorema do Valor Médio para Integrais}
  Seja $f$ uma função {\color{blue} contínua} em $[a, b]$, \pause
  então existe ${\color{blue} c} \in (a, b)$ \pause
  tal que
  \[
    f({\color{blue}c})(b-a)
    \;=\; \int_a^b f(x)\,dx
  \]
  \pause
  ou
  \[
    f({\color{blue}c})
    \;=\; \frac{1}{b-a}\int_a^b f(x)\,dx
  \]
\end{frame}

%------------------------------------------------------------------------------%
\begin{frame}{Valor Médio}
  ${\color{blue}f(c)}$ é o \emph{Valor Médio} de $f$ em $[a, b]$
  que denotamos por ${\color{blue}f_{\text{médio}}}$.
\end{frame}

%------------------------------------------------------------------------------%
\begin{frame}{Valor Médio}
  \centering
  \vspace{\baselineskip}
  \input{fig_valor_medio_media}
\end{frame}

%------------------------------------------------------------------------------%
%------------------------------------------------------------------------------%
%------------------------------------------------------------------------------%
\section{Teorema Fundamental do Cálculo}

%------------------------------------------------------------------------------%
\begin{frame}{Teorema Fundamental do Cálculo -- Parte 1}
  Se $f$ for {\color{blue} contínua} em $[a, b]$,
  \pause
  então a função $F$ definida por
  \[
    F({\color{blue}x})
    \;=\; \int_a^{{\color{blue}x}} f(t)\,dt
    \hspace{15mm} a \leq x \leq b
  \]
  \pause
  é contínua em $[a, b]$,
  \pause
  derivável em $(a, b)$
  \pause
  e
  \[
    \frac{dF}{dx}(x) \;=\; f(x)
  \]
\end{frame}

%------------------------------------------------------------------------------%
\begin{frame}{Em outras palavras}
  Se
  \[
    F({\color{blue}x})
    \;=\; \int_a^{{\color{blue}x}} f(t)\,dt
  \]
  \pause
  então
  \[
    \frac{d}{d{\color{blue}x}}
    \left[\; \int_a^{{\color{blue}x}} f(t)\,dt \;\right]
    = f({\color{blue}x})
  \]
\end{frame}

%------------------------------------------------------------------------------%
\begin{frame}{Teorema Fundamental do Cálculo -- Parte 2}
  Se $f$ for {\color{blue} contínua} em $[a, b]$
  \pause
  então, para {\color{blue}qualquer primitiva} $F$ de $f$
  \[
    \int_a^b f(x) \,dx      \pause
    \;=\; F(b) - F(a)       \pause
    \;=\; F(x) \evalat{a}{b}
  \]
\end{frame}

%------------------------------------------------------------------------------%
\begin{frame}{Usando o TFC para Calcular Áreas}
  Calcular a área sob o gráfico de \(f(x)=x^2\) entre 0 e $b$
  \pause
  \[
    \text{Área}
    \;=\; \int_0^b f(x) \,dx \pause
    \;=\; \int_0^b x^2  \,dx
  \]
  \pause
  Sabemos que a primitiva de $f$ é
  \[
    F(x)
    \;=\; \int x^2  \,dx \pause
    \;=\; \frac{x^{2+1}}{2+1} + C \pause
    \;=\; \frac{x^3}{3} + C
  \]
  \pause
  então
  \[
    \text{Área}
    \;=\; F(x) \evalat{0}{b} \pause
    \;=\; F(b) - F(0)        \pause
    \;=\; \left[ \frac{b^3}{3} + C \right]
      -   \left[ \frac{0^3}{3} + C \right] \pause
    \;=\; \frac{b^3}{3}
  \]
\end{frame}

%------------------------------------------------------------------------------%
%------------------------------------------------------------------------------%
%------------------------------------------------------------------------------%
\section{Demonstrações}

%------------------------------------------------------------------------------%
\begin{frame}{Teorema do Valor Médio para Integrais}
  Queremos mostrar que existe $c \in (a, b)$ tal que
  \[
    f(c) = \frac{1}{b-a}\int_a^b f(x)\,dx
  \]
\end{frame}

%------------------------------------------------------------------------------%
\begin{frame}{Teorema do Valor Médio para Integrais}
  \centering
  \vspace{\baselineskip}
  \input{fig_valor_medio_media}
\end{frame}

%------------------------------------------------------------------------------%
\begin{frame}{Teorema do Valor Médio para Integrais}
\onslide<+->{Teorema de Weierstrass \\
\qquad $f$ é \emph{contínua} em $[a, b]$
então existe um \emph{mínimo} $f(u)$ e um \emph{máximo} $f(v)$}

\vspace{0.5\baselineskip}
\onslide<+->{Como a integral \emph{preserva a desigualdade}}
\begin{align*}
  \onslide<+->{\int_a^b f(u)\,dx \;\leq\; & \int_a^b f(x)\,dx \;\leq\; \int_a^b f(v) \,dx \\[3mm]}
  \onslide<+->{f(u)(b-a) \;\leq\;         & \int_a^b f(x)\, dx \;\leq \;f(v)(b-a)         \\[3mm]}
  \onslide<+->{f(u) \;\leq\;               \frac{1}{b-a} & \int_a^b f(x)\,dx \;\leq\; f(v)}
\end{align*}
\end{frame}

%------------------------------------------------------------------------------%
\begin{frame}{Teorema do Valor Médio para Integrais}
Valor médio
\[
  \fmedio = \frac{1}{b-a} \int_a^b f(x)\,dx
\]
\pause
\[
  f(u) \;\leq\; f_{\text{médio}} \;\leq\; f(v)
\]

\pause
\vspace{\baselineskip}
Teorema do Valor Intermediário \\
\qquad $f$ precisa assumir o valor \(\fmedio\) em algum $c\in[a, b]$
\end{frame}

%------------------------------------------------------------------------------%
%------------------------------------------------------------------------------%
%------------------------------------------------------------------------------%
%------------------------------------------------------------------------------%
\begin{frame}{Teorema Fundamental do Cálculo -- Parte 1}
  Queremos mostrar que, se
  \[
    F(x) = \int_a^x f(t)\,dt
  \]
  então
  \[
    \frac{dF}{dx} = f(x)
  \]
\end{frame}

%------------------------------------------------------------------------------%
\begin{frame}{Teorema Fundamental do Cálculo -- Parte 1}
\onslide<+->{Sejam $x$ e $x+h$ em $(a, b)$}
\begin{align*}
  \onslide<+->{F(x+h) - F(x)
    & = {\color<3>{blue}\int_a^{x+h} f(t)\,dt} - \int_a^x f(t)\,dt                     \\[3mm]}
  \onslide<+->{
    & = {\color<3>{blue}\int_a^x f(t)\,dt + \int_x^{x+h} f(t)\,dt} - \int_a^x f(t)\,dt \\[3mm]}
  \onslide<+->{
    & = \int_x^{x+h} f(t)\,dt}
\end{align*}
\end{frame}

%------------------------------------------------------------------------------%
\begin{frame}{Teorema Fundamental do Cálculo -- Parte 1}
  Considerando $h\neq 0$
  \[
    \frac{F(x+h) - F(x)}{h} = \frac{1}{h} \int_x^{x+h} f(t)\,dt
  \]
  \pause
  Calculando o limite $h\to 0$
  \[
    \lim_{h\to0} \frac{F(x+h) - F(x)}{h} = \frac{dF}{dx}
  \]
  \pause
  Falta mostrar que
  \[
    \lim_{h\to 0} \frac{1}{h} \int_x^{x+h} f(t)\,dt \;\stackrel{?}{=}\; f(x)
  \]
\end{frame}

%------------------------------------------------------------------------------%
\begin{frame}{Teorema Fundamental do Cálculo -- Parte 1}
  Pelo Teorema Valor Médio existe \(c\in[x, x+h]\) tal que
  \[
    \int_x^{x+h} f(t)\,dt = f(c)h
  \]
  \pause
  então
  \[
    \frac{1}{h} \int_x^{x+h} f(t)\,dt = f(c)
  \]
  \pause
  Quando $h\to0$ então $c\to x$, portanto
  \[
    \lim_{h\to 0} \frac{1}{h} \int_x^{x+h} f(t)\,dt = f(x)
  \]
\end{frame}

%------------------------------------------------------------------------------%
%------------------------------------------------------------------------------%
%------------------------------------------------------------------------------%
\begin{frame}{Teorema Fundamental do Cálculo -- Parte 2}
  Queremos mostrar que
  \[
    \int_a^b f(x) \,dx = F(b) - F(a)
  \]
\end{frame}

%------------------------------------------------------------------------------%
\begin{frame}{Teorema Fundamental do Cálculo -- Parte 2}
Teorema Fundamental do Cálculo -- Parte 1
\[
  G(x) = \int_a^x f(x)\,dx
\]
é uma primitiva de $f$

\pause
\vspace{\baselineskip}
Qualquer outra primitiva tem a forma
\[
  F(x) = G(x) + C
\]
\end{frame}

%------------------------------------------------------------------------------%
\begin{frame}{Teorema Fundamental do Cálculo -- Parte 2}
\begin{align*}
  \onslide<+->{F(b) - F(a) & = \big(G(b) + C\big) - \big(G(a) + C\big) \\[5mm]}
  \onslide<+->{            & = G(b) - G(a)                             \\[3mm]}
  \onslide<+->{            & = \int_a^b f(x)\,dx - \int_a^a f(x)\,dx   \\[3mm]}
  \onslide<+->{            & = \int_a^b f(x)\,dx}
\end{align*}
\end{frame}

%------------------------------------------------------------------------------%
\begin{frame}
  \tableofcontents
\end{frame}

%------------------------------------------------------------------------------%
\end{document}
%------------------------------------------------------------------------------%
